\documentclass[a4paper,12pt]{report}
\usepackage{makecell}
\usepackage{placeins}
\usepackage[T5]{fontenc}
\usepackage[american]{babel}
\usepackage[margin=1in]{geometry}
\usepackage{tikz}
\usepackage{amssymb}
\usetikzlibrary{calc}
\usepackage{graphicx}
\usepackage{lmodern}
\usepackage{titlesec}
\titlespacing*{\chapter}{0pt}{0pt}{0pt}
\usepackage{fancyhdr}
\pagestyle{fancy}
\setlength{\headheight}{15pt}
\addtolength{\topmargin}{-2.5pt}
\fancyhead[L]{}
\fancyhead[R]{\textit{\MakeUppercase{\leftmark}}}
\usepackage{amsmath}
\usepackage{array, booktabs, multirow, longtable}
\usepackage{tabularx}
\usepackage{adjustbox}
\usepackage{newunicodechar}
\newunicodechar{≈}{\approx}
\usepackage{listings}
\usepackage{caption}
\usepackage{subcaption}
\usepackage{float}
\usepackage{tocloft}
\usepackage[colorlinks=true, linkcolor=black, citecolor=black, urlcolor=black]{hyperref}
\usepackage{bookmark}
\usepackage{lipsum}
\usepackage{cite}
\usepackage{setspace}
\definecolor{codegreen}{rgb}{0,0.6,0}
\definecolor{codegray}{rgb}{0.5,0.5,0.5}
\definecolor{codepurple}{rgb}{0.58,0,0.82}
\definecolor{backcolour}{rgb}{0.95,0.95,0.92}
\lstdefinestyle{mystyle}{
    backgroundcolor=\color{backcolour},
    commentstyle=\color{codegreen},
    keywordstyle=\color{magenta},
    numberstyle=\tiny\color{codegray},
    stringstyle=\color{codepurple},
    basicstyle=\ttfamily\footnotesize,
    breakatwhitespace=false,
    breaklines=true,
    captionpos=b,
    keepspaces=true,
    numbers=left,
    numbersep=5pt,
    showlines=false,
    showspaces=false,
    showstringspaces=false,
    showtabs=false,
    tabsize=2
}
\lstset{style=mystyle, showlines=false}
% Glossaries are not used in this document, so disable the package to avoid warnings.
% \usepackage[acronym]{glossaries}
% \makeglossaries

\begin{document}

% ===== COVER PAGE =====

\thispagestyle{empty}
\begin{tikzpicture}[remember picture, overlay]
    \draw[line width=2pt]
        ($(current page.north east) - (0.5in, 0.5in)$);
\end{tikzpicture}
\begin{center}
    \adjustbox{max width=\textwidth}{\small \textbf{DEPARTMENT OF INFORMATION AND COMMUNICATION TECHNOLOGY}}\\[2cm]
    \includegraphics[width=0.7\textwidth]{images/vast.png}\\[1.7cm]
    {\Huge \textbf{Report}}\\[1cm]
    \rule{\textwidth}{0.75pt}\\[0.6cm]
    {\LARGE \textbf{Week 4 - Learn to read research papers like a researcher}}\\[0.4cm]
    \rule{\textwidth}{0.75pt}\\[1.2cm]
    {\Large \textbf{By Nguyen Huy Hai, Dang Trung Hieu}}
\end{center}
\vfill
\begin{center}
    % Supervisor above date
    %{\large \textbf{Supervisor: Hoang Do Thanh Tung}} \\ \vspace{0.8cm}
    \textit{Ha Noi, October 2026}
\end{center}

\newpage


% ===== FRONT MATTER =====
% \newpage
% \pagenumbering{arabic}
% \setcounter{page}{1}
\setlength{\parindent}{0pt}
\setlength{\parskip}{3pt}

% \newpage
% \tableofcontents

% \newpage
% \listoffigures
% \addcontentsline{toc}{chapter}{List of Figures}

% \newpage
% \listoftables
% \addcontentsline{toc}{chapter}{List of Tables}

% \newpage
% \chapter*{List of Abbreviations}
% \addcontentsline{toc}{chapter}{List of Abbreviations}
% {
% \renewcommand{\arraystretch}{1.3}
% \begin{longtable}{@{}lp{12cm}@{}}
% \toprule
% Term & Definition \\ 
% \midrule
% WF   & Work Flow \\
% \bottomrule
% \end{longtable}
% }
% \markboth{List of Abbreviations}{}

\newpage
\chapter{Work of this week}

\section{Hieu's work}

\textbf{Do:} Complete a one-page paper dissection for PromptInject, then discuss how goal hijacking differs from harmless task redirection.

\textbf{Deliver:} Paper-reading matrix with 2 entries and 5 open research questions.

\section{Hai's work}

\textbf{Video outline:} Make outline, slide for video LLM vs.\ Agentic AI, and Why Context-Aware Training Matters.

%------------------------------------------------------------
\chapter[R1: Ignore Previous Prompt]{Attack Techniques For Language Models}
 
\section*{Research Problem \& Gap}
 
\textbf{Research Problem:} Large Language Model (LLM) applications frequently construct prompts by inserting untrusted user input directly into pre-written instruction templates by string substitution. This open-ended, unstructured design exposes models to prompt injection attacks, where malicious users embed adversarial instructions into input fields to detail model behavior.
 
\textbf{Research Gap:} While LLMs like GPT-3 had achieved state-of-the-art text processing performance, systematic, quantitative studies examining prompt injection vulnerabilities and automated alignment evaluation frameworks are still difficult to find because the supply is limited or insufficient in academic literature.
 
\section*{Threat Model \& Attack Surface}
 
\textbf{Attacker Profile:} Low-aptitude but ill-intentioned end users operating under a black-box setting, interacting with LLM applications solely through public input fields without access to model weights or gradients.
 
\textbf{Attack Surface:} Unsanitized user input parameters (\texttt{\{user\_input\}}) embedded within prompt templates.
 
\textbf{Attack Objectives:}
\begin{itemize}
  \item \textbf{Goal Hijacking:} Overriding the developer's original prompt instruction to force the model into printing a specific target string or executing an adversarial goal.
  \item \textbf{Prompt Leaking:} Misaligning the model to output private application instructions or secret prompt context.
\end{itemize}
 
\section*{Main Claim \& Hypothesis}
 
\textbf{Core Claim:} GPT-3 and similar state-of-the-art instruction-following models are highly susceptible to handcrafted prompt injection attacks.
 
\textbf{Counter-Intuitive Hypothesis (Inverse Scaling):} More capable, instruction-aligned models (e.g., \texttt{text-davinci-002}) are significantly more vulnerable to prompt injection than weaker models, as their heightened capacity to follow instructions makes them equally adept at following injected malicious commands.
 
\section*{Experimental Setup \& Architecture}
 
\textbf{Framework:} The authors developed \textsc{PromptInject}, a modular framework for mask-based iterative adversarial prompt composition and quantitative robustness scoring.
 
\textbf{Test Dataset:} Evaluated across 35 public base prompts sourced from the official OpenAI Examples directory (e.g., grammar correction, Q\&A, SQL generation, summarization).
 
\textbf{Evaluated Models:} OpenAI API models, primarily \path{text-davinci-002}, alongside \path{text-davinci-001}, \path{text-curie-001}, \path{text-babbage-001}, and \path{text-ada-001}.

\textbf{Experimental Factors:} Systematic variation of attack wording, escape/delimiter characters (\texttt{\textbackslash n----------\textbackslash n}), delimiter length and repetitions, rogue string severity, sampling temperature, top-$p$, penalties, stop sequences, and trailing text after \texttt{\{user\_input\}}. Each setup was executed across 4 runs over all 35 prompts to handle model stochasticity.
 
\section*{Baselines \& Evaluation Metrics}
 
\textbf{Baselines:} Standard prompt responses without attacks, evaluated across different OpenAI GPT-3 models, ranging from smaller, faster models to larger, more capable models.
 
\textbf{Evaluation Metrics:}
\begin{itemize}
  \item \textbf{Goal Hijacking Success Rate (\%):} Percentage of prompts where model output exactly matched the target rogue string (ignoring punctuation/newlines).
  \item \textbf{Prompt Leaking Success Rate (\%):} Percentage of outputs containing the original system prompt instruction.
\end{itemize}
 
\section*{Key Findings \& Defensive Limitations}
 
\begin{itemize}
  \item \textbf{Attack Efficacy:} Reached a success rate of $58.6\% \pm 1.6$ for goal hijacking and $23.6\% \pm 2.7$ for prompt leaking.
  \item \textbf{Structural Delimiters Impact:} Prepending structural delimiters (e.g., \texttt{\textbackslash n----------\textbackslash n}) significantly boosted hijacking success from 43.6\% to 58.6\%.
  \item \textbf{Alignment Backfire on Extreme Threats:} Injected rogue strings containing explicit violence (``Kill all humans'') achieved lower success (49.3\%) than non-harmful strings (70.0\%), indicating partial inhibition from OpenAI's RLHF alignment.
  \item \textbf{Inverse Scaling Confirmed:} \texttt{text-davinci-002} exhibited a 58.6\% attack success rate compared to just 13.8\% in \texttt{text-ada-001}.
  \item \textbf{Defensive Mechanisms:} Hardening techniques like placing static prompt text after \texttt{\{user\_input\}} (reducing success to 51.8\%) and enforcing strict stop sequences (reducing success to 47.5\%) mitigated attacks but failed to prevent them completely.
\end{itemize}
 
%------------------------------------------------------------
\chapter{Matrix comparison between PromptInject and Red Teaming LMs with LMs}
 
\small
\renewcommand{\arraystretch}{1.2}
\setlength{\tabcolsep}{5pt}
\begin{longtable}{@{}>{\raggedright\arraybackslash}p{0.14\textwidth}
                    >{\raggedright\arraybackslash}p{0.40\textwidth}
                    >{\raggedright\arraybackslash}p{0.40\textwidth}@{}}
\toprule
\textbf{Category} & \textbf{PromptInject (R1)} & \textbf{Red Teaming LMs with LMs (R2)} \\
\midrule
\endfirsthead
\toprule
\textbf{Category} & \textbf{PromptInject (R1)} & \textbf{Red Teaming LMs with LMs (R2)} \\
\midrule
\endhead
\bottomrule
\endlastfoot
 
Research Problem &
Untrusted user inputs in prompt templates cause prompt injection attacks that misalign model goals. &
Manual human red teaming is expensive, slow, and unscalable, leaving dangerous model behaviors undiscovered prior to deployment. \\
 
Research Gap &
Quantitative alignment frameworks and empirical benchmarks for prompt injection against commercial LLMs were lacking. &
Scalable automated methods to generate diverse and difficult test inputs eliciting diverse model harms were unstudied. \\
 
Threat Model &
Malicious end users injecting adversarial text via black-box API application inputs. &
Automated Red LM ($p_r$) generating inputs $x$ to probe Target LM ($p_t$) evaluated by Red Team Classifier ($r$) under black-box access. \\
 
Main Claim &
LLM prompt templates are inherently vulnerable to prompt injection; larger instruction-followed models are more vulnerable (inverse scaling). &
Red LMs generate tens of thousands of diverse, difficult failure cases automatically, uncovering risks human red teams miss. \\
 
Experimental Setup &
Modular \textsc{PromptInject} framework testing 35 OpenAI base prompts across multiple attack factors and models. &
Evaluated 280B Dialogue-Prompted Gopher (DPG) probed by 280B and 7B Red LMs across 5 harm domains (offense, leakage, contact info, bias, dialogues). \\
 
Baselines / Target Models &
OpenAI models (\texttt{text-davinci-002}, curie, babbage, ada) under unattacked baseline execution. &
Bot-Adversarial Dialogue (BAD) human-written test cases (2021b) and non-adversarial dialogue generation. \\
 
Evaluation Metrics &
Goal Hijacking Success Rate (\% exact match) \& Prompt Leaking Success Rate (\% substring match). &
Offensiveness Rate (\%), Self-BLEU, $n$-gram entropy, Zipf coefficient, 13-gram verbatim overlap, and regex contact match count. \\
 
Key Findings &
58.6\% hijacking success on davinci-002; delimiters increase attack success; harmful rogue strings lower success; static suffix text mitigates attacks. &
RL Red LMs elicit offensive replies $>$40\% of the time; found 1,709 data leaks, 3,206 phone numbers, and multi-turn offensive escalation. \\
 
Limitations &
Handcrafted attack prompts; no prompt mitigation completely eliminates vulnerability; structural reliance on API constraints. &
Red LMs and classifiers inherit dataset biases (false positives/negatives); prompt engineering sensitivity; offense-defense asymmetry. \\
 
Ethics \& Artifact Availability &
Includes AI X-Risk Analysis (Hendrycks \& Mazeika framework); code open-sourced at GitHub repository. &
Outlines offensive misuse risks vs.\ defensive blue-teaming advantages (rate limits, white-box access, unlikelihood training); research dataset references. \\
 
\end{longtable}
\normalsize
 
%------------------------------------------------------------
\chapter{5 Open Research Questions}
 
\begin{enumerate}
  \item \textbf{Architectural Instruction-Data Separation:} How can transformer architectures be modified natively (e.g., dual-attention streams or instruction/data parameter segregation) to guarantee that user input tensors can never override system instruction contexts?
  \item \textbf{Mitigating Inverse Scaling in Security Robustness:} Why does alignment fine-tuning (e.g., RLHF) increase susceptibility to prompt injection in larger models, and can safety fine-tuning be decoupled from prompt adherence?
  \item \textbf{Automated Synthesis of Indirect Prompt Injection:} How can automated Red LM techniques be adapted to discover multi-modal and indirect prompt injection attacks against tool-integrated AI agents operating in untrusted web environments?
  \item \textbf{Defense Against Conversational Degradation:} Given that offensive and adversarial behaviors escalate over multi-turn dialogues, what lightweight, real-time context-resetting interventions can halt multi-turn escalation without breaking conversation continuity?
  \item \textbf{Unified Evaluation Metrics for Safety-Capability Trade-offs:} How can we establish standardized Pareto frontiers that evaluate a model's safety against automated red-teaming attacks while ensuring the model maintains downstream task performance?
\end{enumerate}
 
%------------------------------------------------------------
\chapter{References}
 
\begin{description}
  \item[R1:] Perez \& Ribeiro (2022), \textit{Ignore Previous Prompt: Attack Techniques For Language Models} --- \url{https://arxiv.org/abs/2211.09527}
  \item[R2:] Perez et al.\ (EMNLP 2022), \textit{Red Teaming Language Models with Language Models} --- \url{https://aclanthology.org/2022.emnlp-main.225/}
\end{description}

\end{document}